\section{Elementary constructions and discrete integration}

\begin{definition}[Measure with prescribed point masses]
\label{def:singleton-measure}
\lean{MeasureTheory.Measure.ofSingletons}
Let $(X,\mathcal A)$ be a measurable space and let
$m:X\to[0,\infty]$. Define the measure
\[
 \mu_m=\sum_{x\in X}m(x)\delta_x,
\]
where a sum of nonnegative extended real numbers is the supremum of
its finite subsums. Thus, for $A\in\mathcal A$,
$\mu_m(A)=\sum_{x\in A}m(x)$.
\end{definition}
\begin{proof}
The value on the empty set is zero. If the measurable sets $A_n$ are
pairwise disjoint, each point occurs in at most one of them, and
\[
 \sum_{x\in\bigcup_n A_n}m(x)
 =\sum_{n=0}^{\infty}\sum_{x\in A_n}m(x).
\]
Indeed, any finite subsum on the left is bounded by a finite sum of
the inner sums on the right. Conversely, any finite choice of finite
subsets of the $A_n$ gives a finite subset of their union, with the
same sum. Taking suprema proves equality. This establishes countable
additivity and hence defines the claimed measure.
\end{proof}

\begin{lemma}[The prescribed mass of a singleton]
\label{lem:singleton-measure-values}
\lean{MeasureTheory.Measure.ofSingletons_singleton}
\uses{def:singleton-measure}
If every singleton in $X$ is measurable, then
$\mu_m(\{x\})=m(x)$ for every $x\in X$.
\end{lemma}
\begin{proof}
In the defining sum, the term indexed by $x$ has value $m(x)$ on
$\{x\}$, while every other Dirac measure assigns zero to that set.
Thus the sum is $m(x)$, including when $m(x)=\infty$.
\end{proof}

\begin{proposition}[Uniqueness from singleton masses]
\label{prop:countable-singleton-measure}
\lean{MeasureTheory.Measure.eq_ofSingletons,MeasureTheory.Measure.existsUnique_singleton}
\uses{def:singleton-measure,lem:singleton-measure-values}
Suppose $X$ is countable and every singleton is measurable. For each
$m:X\to[0,\infty]$, there is exactly one measure $\mu$ with
$\mu(\{x\})=m(x)$ for all $x$; it is $\mu_m$.
\end{proposition}
\begin{proof}
Existence follows from the preceding lemma. Every subset $A$ of $X$
is a countable disjoint union of measurable singletons, hence is
measurable. Countable additivity forces
$\mu(A)=\sum_{x\in A}\mu(\{x\})=\sum_{x\in A}m(x)=\mu_m(A)$.
Thus any measure with the prescribed masses equals $\mu_m$.
\end{proof}

\begin{lemma}[Decomposition of a finitely supported function]
\label{lem:finite-support-indicators}
\lean{MeasureTheory.eq_sum_indicator_of_support_subset}
Let $E$ be a normed abelian group, let $s\subseteq X$ be finite, and
let $g:X\to E$ vanish outside $s$. Then
\[
 g=\sum_{a\in s}\mathbf1_{\{a\}}g(a),
\]
where $\mathbf1_{\{a\}}g(a)$ is $g(a)$ at $a$ and zero elsewhere.
\end{lemma}
\begin{proof}
At a point $x\in s$, exactly the summand indexed by $a=x$ can be
nonzero, and that summand is $g(x)$. At a point outside $s$, every
summand is zero and $g(x)=0$ by hypothesis. This proves equality
pointwise.
\end{proof}

\begin{lemma}[Integrability of finite support]
\label{lem:finite-support-integrability}
\lean{MeasureTheory.integrable_of_support_subset_finset}
\uses{lem:finite-support-indicators}
Let $X$ be a measurable space with measurable singletons, let $\mu$
be a measure on $X$, and let $E$ be a normed abelian group.
If $g:X\to E$ vanishes outside a finite set $s$ and
$\mu(\{a\})<\infty$ for each $a\in s$, then $g$ is integrable.
\end{lemma}
\begin{proof}
For $a\in s$, the function $\mathbf1_{\{a\}}g(a)$ is a measurable
simple function: its only values are zero and $g(a)$, and the
singleton is measurable. It is therefore strongly measurable.
The integral of its norm is
$\mu(\{a\})\lVert g(a)\rVert$, which is finite. Each such function
is consequently integrable. Their finite sum is integrable by the
triangle inequality for the integral of the norm. The preceding
decomposition identifies that sum with $g$.
\end{proof}

\begin{lemma}[Integral of finite support]
\label{lem:finite-support-integral}
\lean{MeasureTheory.integral_eq_sum_of_support_subset_finset}
\uses{lem:finite-support-indicators,lem:finite-support-integrability}
Under the hypotheses of Lemma~\ref{lem:finite-support-integrability},
assume also that $E$ is a real Banach space. Then
\[
 \int_X g\,d\mu=\sum_{a\in s}\mu(\{a\})g(a),
\]
where every coefficient on the right is the finite real mass of the
corresponding singleton.
\end{lemma}
\begin{proof}
The singleton functions in the decomposition are integrable by the
preceding lemma. Linearity of the Bochner integral therefore permits
integration of their finite sum term by term. By the integral formula
for a constant on a measurable set of finite measure,
$\int_X\mathbf1_{\{a\}}g(a)\,d\mu=\mu(\{a\})g(a)$.
Summing these identities proves the formula.
\end{proof}

\begin{definition}[Restriction to a subfield]
\label{def:subfield-restriction}
\lean{ArakelovGeometry.AdelicStructure.fRestriction,ArakelovGeometry.AdelicCurve.fRestriction}
\uses{def:adelic-structure,def:absolute-value-restriction}
Let $S=(K,(\Omega,\mathcal A,\nu),\phi)$ be an adelic curve and
let $F\subseteq K$ be a subfield. Its restriction to $F$ is the
adelic curve
\[
 S|_F=(F,(\Omega,\mathcal A,\nu),
            (\omega\mapsto\phi_\omega|_F)).
\]
\end{definition}
\begin{proof}
Restriction preserves the absolute-value axioms. For $a\in F^\times$,
the logarithmic evaluation for the restricted structure is the same
function as the logarithmic evaluation of $a\in K^\times$ in $S$.
It is therefore measurable and integrable, as required.
\end{proof}

\begin{proposition}[Properness under subfield restriction]
\label{prop:proper-subfield-restriction}
\lean{ArakelovGeometry.AdelicCurve.isProper_fRestriction}
\uses{def:subfield-restriction,def:adelic-structure}
If $S$ is proper, then $S|_F$ is proper.
\end{proposition}
\begin{proof}
For $a\in F^\times$, its image in $K$ is nonzero. The product
formula for this element in $S$ reads
$\int_\Omega\log\phi_\omega(a)\,d\nu(\omega)=0$.
This is exactly its product formula in $S|_F$.
\end{proof}

\begin{definition}[Restriction to a measurable set of parameters]
\label{def:parameter-restriction}
\lean{ArakelovGeometry.AdelicStructure.sRestriction,ArakelovGeometry.AdelicCurve.sRestriction}
\uses{def:adelic-structure}
Let $S=(K,(\Omega,\mathcal A,\nu),\phi)$ be an adelic curve and
let $B\in\mathcal A$. Give $B$ the trace sigma algebra and the
measure $\nu_B$ characterized by
$\nu_B(C)=\nu(C)$ for trace-measurable $C\subseteq B$.
Then
\[
 S|_B=(K,(B,\mathcal A|_B,\nu_B),\phi|_B)
\]
is an adelic curve.
\end{definition}
\begin{proof}
Because $B$ is measurable, a trace-measurable $C\subseteq B$ is
also measurable as a subset of $\Omega$. Thus the displayed formula
defines a measure, and the inclusion $i:B\hookrightarrow\Omega$
sends $\nu_B$ to the restricted measure $\nu|_B$.
For $a\in K^\times$, restriction of the measurable function
$\ell_a(\omega)=\log\phi_\omega(a)$ to $B$ remains measurable.
Moreover,
\[
 \int_B|\ell_a\circ i|\,d\nu_B
 =\int_\Omega\mathbf1_B|\ell_a|\,d\nu
 \leq\int_\Omega|\ell_a|\,d\nu<\infty.
\]
Hence all nonzero logarithmic evaluations remain integrable.
\end{proof}

\begin{definition}[Amalgamation of adelic structures]
\label{def:adelic-amalgamation}
\lean{ArakelovGeometry.AdelicStructure.amalgamation}
\uses{def:adelic-structure}
Let $S_1$ and $S_2$ be adelic curves with common field $K$ and
parameter measure spaces $(\Omega_i,\mathcal A_i,\nu_i)$.
Their amalgamation $S_1\amalg S_2$ has parameter space
$\Omega_1\amalg\Omega_2$, the sigma algebra of sets whose
intersection with each component is measurable, and measure
$(j_1)_*\nu_1+(j_2)_*\nu_2$, where $j_i$ is the component inclusion.
Its absolute value at $j_i(\omega)$ is $\phi_{i,\omega}$.
This defines an adelic curve on $K$.
\end{definition}
\begin{proof}
Each $j_i$ is injective and measurable, and sends measurable sets to
measurable sets by the definition of the disjoint-union sigma algebra.
Thus measurable functions and their integrals on either component
identify with those on its image carrying the pushed-forward measure.
For $a\in K^\times$, the logarithmic evaluation $\ell_a$ restricts
to the measurable function $\ell_{i,a}=\log\phi_{i,\omega}(a)$ on
component $i$. Its inverse image of a Borel set is the union of the
two measurable component inverse images, so $\ell_a$ is measurable.
Finally,
\[
 \int|\ell_a|\,d((j_1)_*\nu_1+(j_2)_*\nu_2)
 =\int|\ell_{1,a}|\,d\nu_1+\int|\ell_{2,a}|\,d\nu_2<\infty.
\]
This proves the required integrability.
\end{proof}

\begin{proposition}[Properness of amalgamation]
\label{prop:proper-amalgamation}
\lean{ArakelovGeometry.AdelicStructure.isProper_amalgamation}
\uses{def:adelic-amalgamation,def:adelic-structure}
If $S_1$ and $S_2$ are proper, then $S_1\amalg S_2$ is proper.
\end{proposition}
\begin{proof}
For $a\in K^\times$, the two component logarithmic evaluations
are integrable. Additivity with respect to the measure and the
change-of-variables formula for the component inclusions give
\[
 \int\log|a|_\omega\,d\nu_{S_1\amalg S_2}(\omega)
 =\int\log\phi_{1,\omega}(a)\,d\nu_1(\omega)
  +\int\log\phi_{2,\omega}(a)\,d\nu_2(\omega).
\]
Both terms vanish by properness, which proves the product formula.
\end{proof}

\begin{definition}[Trivial adelic curve]
\label{def:trivial-adelic-curve}
\lean{ArakelovGeometry.AdelicCurve.trivialAdelicStructure,ArakelovGeometry.AdelicCurve.trivialAdelicCurve}
\uses{def:adelic-structure}
For every field $K$ and every measure space $(\Omega,\mathcal A,\nu)$,
assign the trivial absolute value to each point:
$|0|_\omega=0$ and $|a|_\omega=1$ for $a\ne0$.
These data define an adelic curve.
\end{definition}
\begin{proof}
The trivial absolute value is nonnegative, multiplicative, and vanishes
exactly at zero. If $a+b\ne0$, at least one of $a,b$ is nonzero, so
$|a+b|=1\leq|a|+|b|$; if $a+b=0$ the triangle inequality is immediate.
For every $a\ne0$, the logarithmic evaluation is identically zero.
It is measurable and integrable for every measure $\nu$.
\end{proof}

\begin{proposition}[Properness of the trivial adelic curve]
\label{prop:proper-trivial-adelic-curve}
\uses{def:trivial-adelic-curve,def:adelic-structure}
Every trivial adelic curve is proper.
\end{proposition}
\begin{proof}
For each nonzero field element the logarithmic evaluation is the zero
function. Its integral is zero, which is the product formula.
\end{proof}

\begin{lemma}[Finite measure of a uniformly archimedean locus]
\label{lem:archimedean-locus-finite-measure}
\lean{ArakelovGeometry.AdelicStructure.v_Ωinf_lt_top_of_κ_bddBelow}
\uses{lem:archimedean-bound}
Let $S$ be an adelic curve with archimedean locus $\Omega_\infty$
and exponent $\kappa$. If there is $\varepsilon>0$ such that
$\kappa(\omega)\geq\varepsilon$ for all $\omega\in\Omega_\infty$,
then $\nu(\Omega_\infty)<\infty$.
\end{lemma}
\begin{proof}
The function $\kappa$ is nonnegative and integrable, and
$\Omega_\infty$ is measurable by Lemma~\ref{lem:archimedean-bound}.
The pointwise inequality
$\varepsilon\mathbf1_{\Omega_\infty}\leq\kappa$ gives
\[
 \varepsilon\nu(\Omega_\infty)
 \leq\int_\Omega\kappa\,d\nu<\infty.
\]
Because $\varepsilon$ is strictly positive, the measure of the
archimedean locus must be finite.
\end{proof}

\begin{lemma}[Measurability of the trivial locus]
\label{lem:trivial-locus-measurable}
\lean{ArakelovGeometry.AdelicStructure.measurableSet_ϕ_eq_trivial}
\uses{def:adelic-structure,def:trivial-adelic-curve}
For an adelic structure on a countable field $K$, the set of
parameters whose absolute value is trivial is measurable.
\end{lemma}
\begin{proof}
An absolute value is trivial if and only if it takes the value one
on every element of $K^\times$: its value at zero is already zero.
For $a\in K^\times$, positivity of $\phi_\omega(a)$ gives
\[
 \{\omega:\phi_\omega(a)=1\}
 =\{\omega:\log\phi_\omega(a)=0\}.
\]
This set is measurable because the logarithmic evaluation is
measurable and $\{0\}$ is a Borel set. The trivial locus is the
intersection of these sets over the countable set $K^\times$, hence
is measurable.
\end{proof}
